% nodeflow (on-site) user guide
% build with:  pdflatex nodeflow-user-guide.tex   (run twice for the contents)
% no packages beyond a standard texlive install are needed.

\documentclass[11pt,a4paper]{article}

\usepackage[utf8]{inputenc}
\usepackage[T1]{fontenc}
\usepackage[margin=2.4cm]{geometry}
\usepackage{parskip}
\usepackage{enumitem}
\usepackage{graphicx}
\usepackage{xcolor}
\usepackage{tcolorbox}
\usepackage{titlesec}
\usepackage{fancyhdr}
\usepackage[hidelinks]{hyperref}

\definecolor{leaf}{HTML}{64791F}
\definecolor{olive}{HTML}{5F560F}
\definecolor{ink}{HTML}{232326}
\definecolor{paper}{HTML}{F3F1DF}
\definecolor{warn}{HTML}{A33B1F}

\titleformat{\section}{\Large\bfseries\color{ink}}{\thesection}{0.6em}{}
\titleformat{\subsection}{\large\bfseries\color{olive}}{\thesubsection}{0.6em}{}

\pagestyle{fancy}
\fancyhf{}
\lhead{\small NodeFlow (On-site) user guide}
\rhead{\small\thepage}
\renewcommand{\headrulewidth}{0.4pt}

\newtcolorbox{note}[1][]{colback=paper,colframe=leaf,boxrule=1pt,
  left=8pt,right=8pt,top=6pt,bottom=6pt,#1}
\newtcolorbox{caution}[1][]{colback=white,colframe=warn,boxrule=1pt,
  left=8pt,right=8pt,top=6pt,bottom=6pt,#1}

\newcommand{\nf}{NodeFlow\textsubscript{(On-site)}}
\newcommand{\screen}[2]{%
  \begin{center}\ttfamily\small
  \fcolorbox{ink}{paper}{\parbox{7cm}{\centering #1\\#2}}
  \end{center}}

\begin{document}

\begin{titlepage}
\centering
\vspace*{3cm}
{\Huge\bfseries NodeFlow\textsubscript{(On-site)}\par}
\vspace{0.6cm}
{\Large Sensing System\par}
\vspace{1.4cm}
{\large A guide to building your own soil sensor code\par}
\vspace{0.4cm}
{\large From a blank computer to readings on the screen\par}
\vfill
{\normalsize A research and outreach project of\\
UC Agriculture and Natural Resources and UC Santa Cruz\par}
\vspace{0.6cm}
{\small Version 1, August 2026\par}
\end{titlepage}

\tableofcontents
\newpage

\section{Before you start}

This guide takes you from a computer with nothing installed to a working soil
sensor showing readings on a small screen. You do not need to know how to
program. You will not write a single line of code.

\subsection{What you need}

\begin{itemize}[leftmargin=*]
  \item An Arduino Uno R3 board.
  \item A 16 by 2 LCD keypad shield, the display that presses onto the board.
  \item At least one sensor. The system reads three kinds, listed in
        section~\ref{sec:sensors}.
  \item A USB cable that carries data, not only power. If your computer does
        not notice the board when you plug it in, the cable is the first thing
        to suspect.
  \item A computer running Windows, macOS or Linux.
  \item Power. USB is enough on a bench. For a field, read section~\ref{sec:power}.
\end{itemize}

\subsection{How long it takes}

About twenty minutes the first time, most of which is installing the Arduino
software. After that, generating a new file and uploading it takes two or three
minutes.

\section{Installing the Arduino IDE}

The Arduino IDE is the free program that copies your file onto the board. You
install it once.

\begin{enumerate}[leftmargin=*]
  \item Open \url{https://www.arduino.cc/en/software} in a web browser.
  \item Download the version for your operating system. On Windows choose the
        installer rather than the app from the Microsoft Store, because the
        installer also sets up the drivers the board needs.
  \item Run the file you downloaded and accept the default options. On Windows
        it will ask permission to install device drivers. Say yes to each one.
        On macOS, drag the Arduino icon into your Applications folder.
  \item Open the Arduino IDE once, so it can finish setting itself up.
\end{enumerate}

\begin{note}
On macOS the first time you open it you may see a message saying the app cannot
be checked for malicious software. Open System Settings, go to Privacy and
Security, scroll down, and click Open Anyway.
\end{note}

\subsection{Checking the board is seen}

Plug the Arduino into your computer with the USB cable. A small light on the
board should come on.

In the Arduino IDE:

\begin{enumerate}[leftmargin=*]
  \item Open the \textbf{Tools} menu.
  \item Point at \textbf{Board} and choose \textbf{Arduino Uno}.
  \item Point at \textbf{Port}. You should see an entry that appeared when you
        plugged the board in. On Windows it looks like \texttt{COM3}. On macOS
        it looks like \texttt{/dev/cu.usbmodem} followed by some numbers.
        Choose it.
\end{enumerate}

\begin{caution}
If no new port appears, work through these in order: try a different USB cable,
try a different USB socket, and on Windows check that the driver installed. A
cable that only carries power is by far the most common cause.
\end{caution}

\section{Wiring your sensors}

\subsection{Which pins you can use}

The display shield covers most of the board. What it leaves you is:

\begin{itemize}[leftmargin=*]
  \item \textbf{A1 to A5} for sensors. Five ports.
  \item \textbf{A0} is taken by the buttons on the shield. Do not use it.
  \item \textbf{D2 and D3} are used when a sensor is wired without an adapter,
        described in section~\ref{sec:direct}.
\end{itemize}

\subsection{The sensors this system reads}
\label{sec:sensors}

\subsubsection*{Capacitive soil moisture probe}

A low cost probe, such as the DFRobot SEN0308, with three wires. Red goes to
5V, black goes to GND, and the yellow signal wire goes to one of A1 to A5. It
needs no adapter.

This probe reads \emph{high in dry soil and low in wet soil}. That is a
property of the probe, and the code takes it into account.

\subsubsection*{Irrometer Watermark 200SS}

Measures how hard a plant has to pull to get water out of the soil, which is
called soil water tension and is reported in kPa. A low number means wet soil.
A high number means dry soil.

There are two ways to connect it, and they are genuinely different:

\begin{itemize}[leftmargin=*]
  \item \textbf{Through a 200SS-VA3 adapter.} The adapter does the difficult
        electrical work and hands the board a simple voltage. One adapter
        carries up to three Watermarks and one temperature probe on a single
        port. This is the recommended way, and the way to fit more than five
        sensors.
  \item \textbf{Wired straight to the board.} One sensor only. See
        section~\ref{sec:direct} before you try this.
\end{itemize}

\subsubsection*{Irrometer 200TS soil temperature sensor}

Measures soil temperature. It can sit on its own, or share a 200SS-VA3 adapter
with Watermark sensors.

\subsection{Wiring a sensor without the adapter}
\label{sec:direct}

The form lets you wire a Watermark or a temperature sensor straight to the
board. You will need a 10 kilohm resistor for each one. The sensor goes between
a digital pin and your chosen analog pin, and the resistor goes from that same
analog pin to a second digital pin. The code uses D2 and D3 for this.

\begin{caution}
\textbf{One sensor only.} Wet soil conducts electricity between two bare
sensors, so each one ends up partly reading the other, and over time the metal
inside them corrodes. This is the manufacturer's own warning. If you want more
than one, use a 200SS-VA3 adapter.
\end{caution}

\begin{caution}
\textbf{The combined Watermark and temperature option needs the adapter.} It
exists because the adapter is what joins the two sensors together. The form
does not offer direct wiring for it.
\end{caution}

\section{Powering the board}
\label{sec:power}

On a bench the USB cable powers everything. In a field it does not, so this
section is about what to use instead.

\subsection{How much current the system draws}

\begin{center}
\begin{tabular}{lr}
\hline
\textbf{Part} & \textbf{Roughly} \\
\hline
Arduino Uno, board alone & 45 mA \\
Display backlight & 25 mA \\
Display logic & 2 mA \\
200SS-VA3 adapter & 1.5 mA \\
Capacitive probe & 5 mA \\
\hline
\textbf{Typical total, backlight on} & \textbf{75 mA} \\
\textbf{Typical total, backlight off} & \textbf{50 mA} \\
\hline
\end{tabular}
\end{center}

Only the adapter figure comes from a datasheet. The rest are typical values for
this hardware and vary between boards and shields. Measure yours if runtime
matters: put a multimeter in series with the battery and read the current with
the backlight on and off.

\begin{note}
The backlight is a third of the load. That is why the screen goes dark after
five minutes without a button press, and why any button brings it straight
back. Leaving it on shortens a battery by roughly a third.
\end{note}

\subsection{Two ways to get power in}

\textbf{The barrel socket} is the round connector. It expects 7 to 12 volts and
goes through a regulator on the board that turns the surplus into heat. At 9
volts in and 5 volts out, roughly four tenths of the battery's energy is wasted
warming the regulator. It works, and it is what the low battery warning
watches.

\textbf{The USB socket} expects exactly 5 volts and skips the regulator, so
nothing is wasted. It is the efficient option, but the low battery warning
cannot see anything through it.

\subsection{The options}

\begin{center}
\begin{tabular}{llll}
\hline
\textbf{Battery} & \textbf{Where} & \textbf{Capacity} & \textbf{Rough runtime} \\
\hline
9 V alkaline (PP3) & barrel & 550 mAh & 7 to 11 hours \\
6 AA alkaline & barrel & 2500 mAh & 30 to 50 hours \\
8 AA alkaline & barrel & 2500 mAh & 30 to 50 hours \\
USB power bank & USB & 10000 mAh & 5 to 8 days \\
12 V sealed lead acid & barrel & 7000 mAh & 3 to 6 weeks \\
\hline
\end{tabular}
\end{center}

Runtimes assume the backlight is mostly off. Halve them if it stays on.

\subsubsection*{9 volt alkaline}

The obvious choice, and the worst one. A PP3 holds very little energy for its
size, and the regulator throws away nearly half of what it does hold. Fine for
an afternoon of testing. Do not plan a season around it.

\subsubsection*{AA cells in a holder}

Six AA cells give about 9 volts, which suits the barrel socket. This is the
best simple choice for a few days of logging: AA cells hold four or five times
the energy of a PP3 and cost less.

Eight cells give 12 volts, which lasts no longer in practice, because the extra
three volts are turned into heat rather than reading time.

\begin{caution}
Four AA cells give 6 volts, which is below what the barrel socket needs. The
board will run unevenly or not at all. Use six.
\end{caution}

\subsubsection*{USB power bank}

The most energy for the money, and no regulator waste. Two things to know:

\begin{itemize}[leftmargin=*]
  \item Many power banks switch themselves off when the load is small, and 50
        mA is small. If yours does this, look for one sold as suitable for low
        power devices, or one with a switch that keeps it awake.
  \item \textbf{The low battery warning does not work on USB power.} The board
        cannot tell the difference between a full bank and one that is about to
        cut out. Do not rely on the warning when running this way.
\end{itemize}

\subsubsection*{Sealed lead acid, with or without solar}

For something left in a field. A 12 volt 7 Ah battery runs the system for weeks
and takes a small solar panel and charge controller happily.

Because the regulator wastes the difference between 12 volts and 5, a separate
switching regulator set to 5 volts, feeding the 5V pin rather than the barrel
socket, roughly doubles the runtime. That is a wiring change rather than
something the code knows about.

\subsection{The low battery warning}

The board measures its own supply against a fixed reference inside the chip.
When a battery on the barrel socket runs down far enough that the regulator can
no longer hold 5 volts, the screen shows a warning on the second line.

\begin{caution}
This only works on battery power through the barrel socket. On USB, including
a USB power bank, the supply stays at 5 volts until the moment it disappears,
so no warning can be given. This is a limit of the hardware, not of the code.
\end{caution}

\section{Using the form}

Open the code generator in a web browser:

\begin{center}
\url{https://fernlag.github.io/NodeFlow-Onsite-Sensing-System/}
\end{center}

Once the page has loaded once, it keeps working without an internet connection.

\begin{note}
There is a language button in the top right corner. It switches the whole page
between English and Spanish, and remembers your choice.
\end{note}

\subsection{One block per sensor}

The page starts with one block. Each block describes one sensor plugged into
one port. Use the \textbf{Add another} button if you have more than one sensor.

For each block you answer four questions:

\begin{enumerate}[leftmargin=*]
  \item \textbf{Sensor type.} Which sensor this is.
  \item \textbf{Port.} Which port you plugged it into.
  \item \textbf{What you want to see on the screen.} The measurement, described
        in section~\ref{sec:measurements}.
  \item \textbf{How it should appear.} A number, a percentage, a bar, or words
        describing the state.
\end{enumerate}

\subsection{The settings underneath}

Below those four questions is a section of settings. It only shows the ones
that matter for the choices you made. If you pick a measurement that needs
nothing configured, the section disappears entirely. This is deliberate: a
setting you should not touch is not shown.

Each setting has a small \textbf{i} button beside it. Press it, or move onto it
with the Tab key, and it explains what the number means and where to get it.

\begin{caution}
The default values are starting points, not measurements of your soil. A
default is a guess. The two calibration values for a capacitive probe in
particular must be measured, and section~\ref{sec:calibration} explains how.
\end{caution}

\subsection{Generating the file}

Press the large \textbf{Generate} button at the bottom. A window appears asking
for your name, country and email, and for a name for the file.

This information goes to the research project. The consent form is linked in
that window, and you have to tick the box before the button becomes available.
Nothing is sent until you press confirm. The privacy policy on the site
explains exactly what is recorded and how to have it removed.

If you tick the date and time option, the file name gets a stamp on the end, so
each download sits alongside the last one instead of being numbered by your
browser.

The file downloads to wherever your browser puts downloads. It ends in
\texttt{.ino}, which is the Arduino file type.

\section{Putting the file on the board}

\begin{enumerate}[leftmargin=*]
  \item Find the downloaded file and open it. The Arduino IDE starts.
  \item The IDE offers to move the file into a folder of the same name. Say
        yes. That is simply how Arduino files are stored.
  \item Check \textbf{Tools} $\rightarrow$ \textbf{Board} still says Arduino
        Uno, and \textbf{Tools} $\rightarrow$ \textbf{Port} still points at
        your board.
  \item Press the arrow button in the top left to upload. Lights on the board
        flicker while it copies.
  \item When the bar at the bottom says \textbf{Done uploading}, the display
        starts showing readings.
\end{enumerate}

\begin{note}
Nothing needs installing for the code to work. It uses only the display library
that comes with the Arduino IDE.
\end{note}

\section{Reading the screen}

The top line shows the port, a short sensor name, and the value.

\screen{A2 WM200SS  34kPa}{OK\ \ \ irrigate}

The bottom line shows whatever you chose: a bar, a percentage, or words. Press
the \textbf{left} and \textbf{right} buttons on the shield to step through your
sensors.

The backlight switches off after five minutes without a button press, to save
power. Any button wakes it.

\subsection{Which way the numbers run}
\label{sec:direction}

This is worth reading once, because it prevents a common misunderstanding.

\textbf{Every percentage runs from dry to wet.} Zero means dry. One hundred
means wet. This is true for every percentage the system shows, without
exception.

\textbf{Tension in kPa runs the other way,} because that is what tension is.
Zero kPa is soaking wet soil. A high number is dry soil. So that this cannot be
confused, the second line of the display always spells out the state next to
the number, as in the example above.

\textbf{Raw readings report the sensor itself.} A capacitive probe reads high
in dry soil, so a raw reading from it goes up as the soil dries. If you want a
number that goes up as soil gets wetter, choose the percentage rather than the
raw value.

\section{The measurements you can choose}
\label{sec:measurements}

\begin{description}[leftmargin=1cm,style=nextline]
  \item[Raw value] Exactly what the sensor put on the pin, with nothing done to
        it. Useful for calibration and for checking a sensor is alive.
  \item[Raw value as a percentage] The raw reading placed on a scale between
        your dry and wet calibration points. Zero is dry, one hundred is wet.
  \item[Management thresholds] Three plain states rather than a number: too
        dry, about right, or saturated. For a Watermark you choose your soil
        texture and the two tension thresholds fill themselves in.
  \item[Tension in kPa] Soil water tension. Low is wet, high is dry.
  \item[Temperature] Soil temperature, in Fahrenheit or Celsius.
  \item[Wetting front] Whether irrigation water has reached a given depth yet.
        Needs two probes at different depths, so it needs two blocks.
  \item[Vertical flow rate] How fast water travelled between the two depths,
        in centimetres per hour. Also needs two probes.
  \item[Volumetric soil moisture] How much of the soil's volume is water,
        as a percentage.
  \item[Total available water] How much of the water the soil can hold is
        actually available to a crop. Zero at the wilting point, one hundred at
        field capacity.
\end{description}

\section{Calibrating a capacitive probe}
\label{sec:calibration}

A capacitive probe needs two numbers before a percentage means anything: what
it reads in air, and what it reads in water. They differ between probes, so
measure yours.

\begin{enumerate}[leftmargin=*]
  \item In the form, choose \textbf{Raw value} as the measurement for that
        sensor, and \textbf{Raw value} as the display. Generate the file and
        upload it.
  \item Hold the probe in open air, away from your hand. Write down the number
        on the screen. It will be somewhere near 590.
  \item Stand the probe in a glass of water, up to but not past the line marked
        on it. Write down that number. It will be somewhere near 280.
  \item Go back to the form, choose the measurement you actually want, and put
        your two numbers into the air value and water value boxes.
  \item Generate the file again and upload it.
\end{enumerate}

\begin{caution}
Do not put the probe in water past the marked line. The electronics at the top
are not waterproof.
\end{caution}

\section{When something looks wrong}

\begin{description}[leftmargin=1cm,style=nextline]
  \item[The screen is blank, or shows a row of solid blocks]
        Turn the small screw on the shield, which is the contrast control. A
        new shield is often set to one extreme or the other.
  \item[A port says No Sensor]
        Nothing is detected on that pin. Check both ends of the wiring. If the
        sensor goes through a 200SS-VA3 adapter, give it time: the adapter
        updates each sensor once every five minutes, so a reading can be up to
        five minutes old.
  \item[The percentage seems backwards]
        Your air and water calibration numbers are probably the wrong way
        round. See section~\ref{sec:calibration}.
  \item[A battery warning appears while on USB]
        Ignore it. The check only means something on battery power through the
        barrel socket. See section~\ref{sec:power}.
  \item[The board keeps switching off in the field]
        If it runs from a USB power bank, the bank is probably shutting down
        because the load is too small for it to notice. See
        section~\ref{sec:power}.
  \item[The numbers do not match my soil]
        The calibration values and the soil thresholds belong to your field,
        not to us. Change them in the form and generate the file again.
  \item[The upload fails]
        Check the port is still selected. Unplug the board, plug it back in,
        select the port again and retry. If it still fails, close the Serial
        Monitor window if you have it open, as it can hold the port.
\end{description}

\section{Working without an internet connection}

The page works offline once it has been opened online at least once. Fill in
the form and download files with no connection at all.

Clearing your browser's history or site data for this page removes the offline
copy, along with your saved details. Open the page once more while online and
it comes back.

\section{Getting help}

The contact address is at the bottom of every page on the site. Tell us which
sensor, which port, which measurement, and what the screen shows. If you can,
include the file name you generated: it lets us find your exact configuration.

\vfill
\begin{center}
\small The generated code is provided as is, without warranty. Before you
irrigate on the strength of a reading, check that it matches what the soil is
actually doing: dig, or compare it against an instrument you already trust.
\end{center}

\end{document}
